\origpage{61}{GRAPH MINORS}
\noindent suitably chosen. This is a contradiction, and so $C$ has no such rooted minor.
By induction, (5.1) is true.
\stmt
(5.2). \textit{If $X_0$ contains a vertex of $C$, then $C$ has path-width at most
$q(\gamma, 3\Omega+6)$.}
\stmt
In this case $C$ contains no vertices of $\bigcup(X_j\colon j\in\{1,..., N\}-J)$ since $X$
separates $X_0$ and this set. Let $P$ be a minimal path of $B_0$ from $V(C)$ to $\rho(G)$.
Let $v$ be the vertex of $C$ on $P$. Root $C$ at $v$. Suppose that $C$ has a rooted
minor isomorphic to $Q(\gamma, 3\Omega+6)$. Then there exist $X_0',..., X_{\Omega+1}'\subseteq V(C)$,
disjoint, such that $C\,|\,X_0'$, rooted at $v$, has a rooted minor isomorphic to $H$,
and for $1\le i\le\Omega+1$, $C\,|\,X_i'$ has a minor isomorphic to $\mathbf{Y}_{\gamma-1}$. But then the
sets
\begin{equation*}
X_0'\cup V(P), X_1',..., X_{\Omega+1}'\qquad\text{and}\qquad X_j(j\in\{1,..., N\}-J)
\end{equation*}
satisfy (X1) and (X2), and there are $\Omega+1+N-|J|>N$ of them. This
contradicts the maximality of $N$, and hence proves (5.2). This completes the
argument for Part 3, and so proves (3.7).

\sectionhead{\textsc{References}}

{
\newcommand{\refauth}{N. \textsc{Robertson and} P. D. \textsc{Seymour}}
\begin{list}{}{\leftmargin=16pt\labelwidth=12pt\labelsep=4pt\itemsep=0pt\parsep=0pt}
\item[1.] \refauth, Graph minors. II. Algorithmic aspects of tree-width,
submitted.
\item[2.] \refauth, Graph minors. III. Planar tree-width, \textit{J.
Combinatorial Theory} (Ser. B), to appear.
\item[3.] \refauth, Graph minors. IV. Tree-width and well-quasi-ordering, submitted.
\item[4.] \refauth, Graph minors. V. Excluding a planar graph,
submitted.
\end{list}}
